\documentclass[10pt]{article}
\usepackage[margin=1in]{geometry}
\usepackage[T1]{fontenc}
\usepackage{lmodern}
\usepackage{amsmath,booktabs}
\usepackage[round]{natbib}
\usepackage{xcolor}
\usepackage{hyperref}

% revision marking from the first round
\definecolor{revblue}{RGB}{0,0,200}
\newcommand{\rev}[1]{{\color{revblue}#1}}
% second round
\definecolor{neworange}{RGB}{214,96,0}
\newcommand{\new}[1]{{\color{neworange}#1}}

\title{Retrieval-Augmented Transformers for Tabular Prediction}
\author{Anonymous}
\date{}

\begin{document}
\maketitle

\section{Introduction}
We pair a transformer encoder with a nearest-neighbor retrieval step over the
training table. \rev{Retrieval uses the same embedding space as the encoder, so
no second model is trained.} The encoder follows residual designs
\citep{he2016resnet}.

\section{Method}
Each query row retrieves its $k=32$ nearest training rows. \new{We set $k$ by
validation accuracy on a held-out tenth of the training rows.} The model is
trained with Adam \citep{kingma2015adam}.

\section{Results}
Table~\ref{tab:main} reports accuracy on the benchmark.
\begin{table}[h]
  \caption{\textbf{Accuracy on the benchmark.} Mean over five seeds.}\label{tab:main}
  \centering
  \begin{tabular}{@{}lc@{}}\toprule Model & Accuracy \\\midrule Gradient boosting & .842 \\ Ours & .861 \\\bottomrule\end{tabular}
\end{table}

\section{Conclusion}
Retrieval improves tabular transformers at no extra training cost.

\bibliographystyle{plainnat}
\bibliography{refs}
\end{document}
